\documentclass[10pt,twocolumn]{article}
\usepackage[T1]{fontenc}
\usepackage[utf8]{inputenc}
\usepackage{lmodern}
\usepackage[margin=1.8cm,top=2.4cm,bottom=2cm,columnsep=0.6cm]{geometry}
\usepackage{fancyhdr}
\usepackage{titlesec}
\usepackage{booktabs}
\usepackage{amsmath,amssymb,amsfonts}
\usepackage{microtype}
\usepackage{xcolor}
\usepackage{mdframed}
\usepackage{parskip}
\usepackage{enumitem}
\usepackage{hyperref}

\definecolor{rulered}{RGB}{180,30,30}
\definecolor{boxgray}{RGB}{245,245,245}
\definecolor{keywordgray}{RGB}{100,100,100}

\pagestyle{fancy}
\fancyhf{}
\fancyhead[L]{\textsc{\small Claude Mag}}
\fancyhead[C]{\small\textit{Machine Learning Research Digest}}
\fancyhead[R]{\small 2026-09-23}
\fancyfoot[C]{\small\thepage}
\renewcommand{\headrulewidth}{0.8pt}

\titleformat{\section}{\normalsize\bfseries\color{rulered}}{}{0pt}{}%
  [\vspace{-4pt}\color{rulered}\rule{\columnwidth}{0.4pt}]
\titlespacing{\section}{0pt}{8pt}{4pt}

\hypersetup{colorlinks=true,urlcolor=rulered,linkcolor=black}

\begin{document}
\setlength{\parindent}{0pt}
\setlength{\parskip}{4pt}

{\small\color{keywordgray}\textbf{Keywords:}
  sim-to-real transfer \textbullet{}
  contact-rich manipulation \textbullet{}
  reinforcement learning \textbullet{}
  force feedback}\par\vspace{4pt}

{\LARGE\bfseries\raggedright
  InsertAnything}\par\vspace{4pt}

{\large\itshape\raggedright
  A Single Simulation-Trained Policy Achieves 95\% Success Across Eight
  Real-World Precision Insertion Tasks at 0.02\,mm Clearance}\par\vspace{6pt}

{\small\textbf{By} [Authors, arXiv:2609.24511]}\par\vspace{2pt}

{\small\color{keywordgray}arXiv:2609.24511 \textbullet{} September 2026}
\par\vspace{2pt}\noindent{\color{rulered}\rule{\columnwidth}{0.6pt}}\vspace{6pt}

InsertAnything trains a reinforcement-learning insertion policy entirely
in simulation using only target poses and compact 3-axis fingertip forces,
then deploys it directly to real robots---achieving 95.0\% success across
eight unseen geometries at sub-millimeter clearances and the first perfect
score on ManipulationNet's peg-in-hole benchmark.

\begin{mdframed}[backgroundcolor=boxgray,linecolor=rulered,linewidth=1pt,
    innerleftmargin=6pt,innerrightmargin=6pt,
    innertopmargin=6pt,innerbottommargin=6pt]
\textbf{\small AT A GLANCE}\par\vspace{4pt}
\begin{description}[leftmargin=0pt,labelsep=4pt,itemsep=2pt,parsep=0pt,
    font=\small\bfseries]
  \item[Problem:] Precision peg-in-hole insertion at sub-mm clearances
    requires geometry-specific demonstrations; policies do not transfer
    across shapes.
  \item[Why it matters:] Industrial assembly requires a single robot to
    insert hundreds of distinct parts; retraining per geometry is
    impractical.
  \item[Method:] RL in simulation with decoupled gated reward, 3D
    fingertip force feedback, force smoothing, and state-independent
    action standard deviations.
  \item[Result:] 95.0\% success across 8 real insertion tasks; 20/20 on
    ManipulationNet; min clearance 0.02\,mm.
  \item[Evidence:] Cross-geometry and cross-clearance evaluation;
    4-axis ablation study confirming each design choice.
\end{description}
\end{mdframed}

\section{The Big Picture}

Precision insertion---placing a peg or bolt into a hole with clearances
as small as 0.02\,mm---is an archetypal industrial assembly skill. It
requires tolerating position estimation errors, searching for alignment
under contact, and recovering from jams. These challenges compound
dramatically across geometries: the contact mechanics for a hexagonal
peg differ from a round one, and the jam-recovery strategy that works
for one fails for the other.

The standard robotics pipeline collects new demonstrations or builds a
new controller for each geometry. InsertAnything takes the opposite
approach: train a single RL policy in simulation on \emph{one} geometry
(hexagonal), using only a minimal sensor suite---target poses plus compact
3D fingertip forces---and show that the resulting policy generalizes to
eight qualitatively different real geometries it has never seen.

The key insight is that the \emph{skills} required for insertion---lateral
alignment search, jam detection via force, and downward insertion progress
once aligned---are geometry-agnostic. The policy must learn these skills,
not geometry-specific trajectories.

\section{Why Existing Approaches Fall Short}

\textbf{Imitation learning from demonstrations} binds generalization to
the demonstrated geometries; re-collection is required for each new part.

\textbf{Model-based contact planning} requires accurate part geometry and
contact-mode enumeration, intractable for arbitrary shapes at runtime.

\textbf{Vision-based RL} adds perception errors on top of control errors
and requires careful camera calibration; contact state is ambiguous from
images.

\textbf{Rich tactile sensors} (GelSight, TacTip) provide contact images
but introduce high-dimensional observation spaces and complex sim-to-real
transfer for the sensor model itself.

InsertAnything uses the simplest possible contact signal---3D force
vectors per fingertip---which is easy to simulate and calibrate to real
hardware, while being sufficient to carry the alignment and jam-recovery
information needed for precise insertion.

\section{Core Mechanism}

\textbf{Observation Space.} Two compact inputs:
\begin{enumerate}[itemsep=2pt,parsep=0pt]
  \item A 6-DoF target pose for the end-effector
  \item Smoothed 3D fingertip force readings $\tilde{\mathbf{f}}_t$
\end{enumerate}
No camera input. The force signal provides the sole contact feedback.

\textbf{Decoupled Gated Reward.} The reward is gated by an alignment
threshold $\epsilon_{\mathrm{align}}$:
\[
r_t = \begin{cases}
r_{\mathrm{align}}(s_t) & \text{if lateral error} > \epsilon_{\mathrm{align}} \\
r_{\mathrm{insert}}(s_t) & \text{otherwise}
\end{cases}
\]
$r_{\mathrm{align}}$ rewards lateral error reduction; $r_{\mathrm{insert}}$
rewards downward depth progress. The gate prevents premature insertion
reward before alignment, eliminating the jam-inducing failure mode.

\textbf{Force-Signal Smoothing.} A sliding-window average reduces noise
without lag:
\[
\tilde{\mathbf{f}}_t = \frac{1}{W}\sum_{w=0}^{W-1}\mathbf{f}_{t-w}
\]

\textbf{State-Independent Action Standard Deviations.} The stochastic
policy uses learned but state-independent $\boldsymbol{\sigma}$, preventing
entropy collapse during the precision insertion phase and preserving
exploratory behavior for jam recovery.

\section{Technical Formulation}

Policy:
\[
\pi_\theta:
(\mathbf{p},\,\tilde{\mathbf{f}})\;\mapsto\;
\mathcal{N}\!\left(\boldsymbol{\mu}_\theta(\mathbf{p},\tilde{\mathbf{f}}),\,
\mathrm{diag}(\boldsymbol{\sigma}^2)\right)
\]
where $\mathbf{p}\in\mathbb{R}^7$ is the target pose and
$\tilde{\mathbf{f}}\in\mathbb{R}^{3n_f}$ the smoothed forces for $n_f$
fingertips.

Alignment reward (lateral distance to hole axis):
\[
r_{\mathrm{align}} = -\!\sqrt{(x-x^*)^2 + (y-y^*)^2}
\]

Insertion reward (depth progress):
\[
r_{\mathrm{insert}} = z^* - z_t
\]

Full episode return:
\[
R = \sum_t \gamma^t r_t
  + R_{\mathrm{success}}\cdot\mathbf{1}[\mathrm{depth}\ge d_{\mathrm{target}}]
\]

\section{Learning or Inference Procedure}

\begin{enumerate}[itemsep=2pt,parsep=0pt]
  \item Randomise hole position ($\pm5$\,mm), peg orientation ($\pm15°$),
    friction in MuJoCo; train PPO on hexagonal insertion only.
  \item Calibrate force scaling (simulation to real sensor).
  \item Deploy on real robot---zero demonstrations, zero fine-tuning.
  \item Evaluate on 8 novel real geometries (round, square, triangular,
    star-shaped, L-slot, and more).
\end{enumerate}

\section{What the Guarantee Says}

The decoupled gated reward provides a structural inductive bias: the
policy \emph{cannot} earn insertion reward without first satisfying
the alignment criterion. This eliminates the reward-hacking failure mode
where policies achieve partial insertion depth by wedging the peg
in a jammed position, which maximises depth reward but causes part damage.

\section{Experimental Findings}

\begin{itemize}[itemsep=2pt,parsep=0pt]
  \item \textbf{Main result:} 95.0\% average success over 8 real tasks;
    minimum clearance 0.02\,mm
  \item \textbf{ManipulationNet:} First-ever 20/20 perfect score on the
    peg-in-hole benchmark under the Human-in-the-Loop protocol with fully
    autonomous insertion motions
  \item \textbf{Cross-clearance:} Policy trained at 0.1\,mm generalises
    to 0.05\,mm with modest success degradation
  \item \textbf{Force reduction:} Peak contact force $\approx$30\% lower
    than baseline without force feedback, reducing part wear
\end{itemize}

\section{Ablations and Interpretation}

\textbf{No force feedback (pose-only):} Success drops to $\approx$40\%;
the policy cannot recover from jams without force information.

\textbf{No decoupled gate (linear sum reward):} Success drops to
$\approx$65\%; premature insertion attempts cause hard jams.

\textbf{State-dependent $\sigma$:} Success drops to $\approx$70\%; the
policy becomes near-deterministic during insertion, losing jam-recovery
exploration.

\textbf{No force smoothing:} Success drops to $\approx$80\%; noisy force
signals cause erratic recovery oscillations that destabilise alignment.

\section*{Reference}

\textbf{InsertAnything: Generalizable Contact-Rich Precision Insertion
from Simulation to Reality}.
arXiv:2609.24511, September 2026.\\
\url{https://arxiv.org/abs/2609.24511}

\end{document}
